\chapter{Introduction}

\begin{quote}
\enquote{Qu'est-ce qu'une bibliothèque numérique ? Il est permis de dire que lors de ce passage du monde physique au monde numérique, tout change et rien ne change. Tout change puisque la puissance et la rapidité de l'outil numérique bouleversent les modes d'accès ; mais rien ne change, car la capacité d'une bibliothèque à répondre à la demande des lecteurs -- qu'elle soit d'information, de savoir, de réflexion ou de distraction -- continue de dépendre de la qualité de son classement et de ses outils descriptifs. Hier comme aujourd'hui, ils sont le fruit du travail et de la compétence des professionnels qui en ont la charge[...]. Le bibliothécaire demeure l'acteur-clé ; l'accès qui dépendait hier du bon catalogage continue, à l'heure du web sémantique et des données liées, de reposer sur son travail de structuration, de description et d'indexation des documents.}
\end{quote}
\medskip

\lettrine{D}ans la préface du\textit {Manuel de constitution de bibliothèques numériques}, Bruno Ovry-Lavollée a ces mots très jutes pour décrire une bibliothèque numérique et son impact sur les conditions de travail, demeure aujourd'hui encore d'une profonde justesse. Le travail du bibliothécaire est la clé de voûte sur laquelle repose tout le projet d'une bibliothèque numérique. Pour lui, le concept d'une bibliothèque numérique est abstrait et vaste, rendant difficile à quiconque souhaitant la définir d'en déterminer tous les aspects. Par ailleurs, il fait remarquer à juste titre que le titre de l'ouvrage mentionne au pluriel le terme de bibliothèque numérique, laissant sous-entendre qu'il n'existerait pas qu'une seule bibliothèque numérique, mais bien plusieurs. En réalité, elle représenterait un organisme dans lequel doit s'insérer les nouvelles bibliothèques numériques afin de pouvoir communiquer entre elles et évoluer en s'appuyant mutuellement dessus.

\enquote{De même que des cellules forment un organe, ces organes un corps, et ces corps une famille, une meute ou une société, les bibliothèques numériques doivent organiser leur complémentarité et leurs interactions. La création d'une bibliothèque numérique ne saurait être pensée qu'en fonction de son insertion parmi les autres et des relations avec elles}.
\medskip

L'objectif de ce mémoire est de montrer l'évolution de ce qu'est une bibliothèque numérique de ce qu'elle représente pour les institutions patrimoniales. Au cours de ces dernières décennies, les bibliothèques numériques représentent un enjeu important en terme de visibilité. En effet, le public touché en ligne n'est pas forcément le même que celui qui consulte habituellement les collections dans la bibliothèque. Aujourd'hui, la bibliothèque numérique est passée d'un simple catalogue de PDF à un site web interactif présentant différents services cherchant à attirer le lecteur vers ses collections en les valorisant à travers des expositions virtuelles, des billets de blog, des lettres d'actualités. La bibliothèque numérique est considérée comme une vitrine de l'institution patrimoniale à laquelle elle se trouve rattacher. Il s'agit d'une extension de la bibliothèque physique en ligne qui doit donc respecter des normes et standards bibliographiques afin que celle-ci puisse communiquer avec les Systèmes Intégrés de Gestion de Bibliothèque (SIGB) et avec les autres bibliothèques numériques. Face à cette masse d'information submergeant les utilisateurs, la visibilité des bibliothèques numériques demeure un véritable enjeu difficile à surmonter. L'une des solutions est l'interopérabilité plus la bibliothèque respecte les normes en vigueur, utilise les mêmes outils permet de s'appuyer sur la communauté qui se crée in fine et c'est grâce à celle-ci qu'elle est capable d'améliorer la visibilité de ses collections.

\bigskip
Ce travail repose sur une démarche empirique et réflexive, nourrie d'un stage de quatre mois au sein du Département du patrimoine, de la conservation et de la numérisation (DPCN) de la Bibliothèque interuniversitaire Cujas (BIU Cujas). Le but est de permettre à la bibliothèque de disposer d'une première étude pour réinformatiser leur bibliothèque numérique.

Ce stage s'est articulé autour de deux axes principaux. Le premier visait à observer et documenter l'existant, par un état des lieux fonctionnel de la bibliothèque numérique CujasNum. Cet état des lieux prend en compte les points de vue des gestionnaires et des usagers, en examinant plusieurs aspects : la plateforme logicielle, l'architecture technique, l'ergonomie, l'accessibilité, etc. Le second axe consistait à identifier les besoins fonctionnels de la future bibliothèque numérique. Pour cela, une enquête accompagnée d'entretiens auprès d'un échantillon de personnes -- les usagers et le personnel de la BIU Cujas -- afin de mieux cerner leurs attentes concernant la future plateforme. A l'issue de ces entretiens, une synthèse a été réalisé mettant en avant les différentes fonctionnalités que le public attendait. Cette base a servie de ligne de conduite pour réaliser le benchmark auprès des autres bibliothèques numériques.

En parallèle de cette étude menée sur l'existant, il convenait d'analyser le prototype Codex fraîchement présenté au DPCN. Cette démarche excluait du benchmark l'étude comparative des logiciels techniques, tels qu'Omeka S ou Limb Gallery. La création de ce prototype répond en réalité à une urgence à laquelle la BIU Cujas risque de faire face. Il s'agit de la possibilité d'une fermeture à tout moment du logiciel sous lequel fonctionne CujasNum. En effet, l'infrastructure actuelle de CujasNum sous XTF, devenue obsolète et non mise à jour depuis 2018, fait peser sur l'établissement un risque de fermeture brutale de la bibliothèque numérique pour des raisons de sécurité informatique. Ne souhaitant pas voir ce service interrompu, il a donc fallu développer une solution de repli. Face à cet impératif de continuité de service, le prototype permet donc un basculement de la bibliothèque numérique afin de permettre aux usagers de pouvoir continuer d'y avoir sans être dépaysée sur son fonctionnement. Le prototype repose sur la version actuelle de la bibliothèque numérique. Codex est développé sous Python-Django, conçu en interne comme une alternative souveraine et immédiate pour parer à l'éventualité d'une fermeture forcée de la plateforme historique.

\medskip
Ce passage d'un système à l'autre soulève des questions profondes qui dépassent le simple cadre de l'outil informatique. Dès lors, une question centrale guidera notre réflexion : dans quelle mesure la refonte de la bibliothèque numérique Cujas impose-t-elle de repenser son architecture logicielle et l'organisation de ses pratiques professionnelles afin de s'affranchir de son isolement et de sa dette technique ?
Ce mémoire s’articulera autour de trois axes principaux, afin de répondre à cette problématique.

Tout d'abord, il conviendra de dresser un diagnostic historique et technique de l’accès au patrimoine numérique de la bibliothèque Cujas. Nous analyserons comment l’institution s’est affirmée comme une référence de la recherche juridique avant de retracer la création de CujasNum pour mettre en lumière la dette technique et l'obsolescence ergonomique qui menacent aujourd’hui la pérennité de sa diffusion.

Puis, nous combinerons à l'historique de création  de Cujas l'expérience de terrain à travers la collecte des besoins des usagers aussi du côté des lecteurs que du côté des personnels de Cujas. Nous essayerons de voir comment nous pourrons insérer ces besoins au sein des réseaux documentaires interconnectés. Il s'agira d'étudier comment les principes du Web sémantique et de la Transition bibliographique permettent de briser les silos documentaires, tout en mesurant, à l'aide d'une enquête empirique, les besoins et les usages réels de la communauté des chercheurs, des étudiants et des professionnels du droit.

Enfin, après avoir identifié les différents besoins fonctionnels, nous aborderons la conception des services et d’accompagnement du changement dans un système en mutation. En effet, le monde des bibliothèques est plein changement donc il est important qu'à travers l’étude du prototype Codex et d'un benchmark ciblé de le prendre en considération afin de répondre au mieux aux besoins. Nous traduirons ces attentes en spécifications fonctionnelles (visionneuse, espace personnel, éditorialisation) avant d’interroger les mutations organisationnelles et professionnelles indispensables pour pérenniser cette transition.